\chapter{完备空间：贝尔、阿斯科利、斯通–魏尔斯特拉斯}
\label{ch:b3:complete}

完备性——每个柯西列都收敛——是
使分析能够\emph{产生}对象的性质：
压缩映射的不动点、级数的和，以及
以极限得到的方程解。本章汇集度量理论中三大存在性
工具。\emph{贝尔定理}表明
完备空间不能是可数个可忽略部分的并，
并凭空造出对象（处处不可导的
连续函数！），靠的只是纯粹的基数式推理。
\emph{阿尔泽拉–阿斯科利}刻画了
$\mathcal C(K)$ 的紧子集，是分析中紧性的主力工具——
周末问题借助它证明皮亚诺关于
微分方程的存在性定理。\emph{斯通–魏尔斯特拉斯}表明
多项式及其他许多函数在 $\mathcal C(K)$ 中稠密：
逼近化为代数验证。沿途
构造完备化，并证明
一致连续映射的延拓定理，它们是
\cref{ch:b3:lp,ch:b3:hilbert,ch:b3:fouriertransform}中的日常工具。

\section{完备空间、完备化与延拓}

\begin{definition}\label{def:b3:complete:complete}
若度量空间中每个柯西列都收敛，则称其为\emph{完备度量空间}\index{完备度量空间}
（本科二年级已知：$\R^n$ 完备；
$\mathcal C(\intcc01)$ 配上 $\norm\cdot_\infty$ 后完备）。
完备空间的闭子集是完备的；
任一度量空间的完备子集都是闭的。
\end{definition}

\begin{proof}
对于两个结论：闭集 $F$ 中的柯西列
在 $X$ 中收敛，其极限是 $F$ 的附着点，因而属于 $F$；
在 $X$ 中收敛、取自完备集 $A$ 的序列是柯西列，
从而在 $A$ 中收敛，且极限唯一。
\end{proof}

\begin{theorem}[一致连续映射的延拓]
\label{thm:b3:complete:extension}
设 $D \subseteq X$ 为稠密子集，$Y$ 为完备空间，且 $f \colon D \to
Y$ 为一致连续映射。则 $f$ 可\emph{唯一}地延拓为
连续映射 $\bar f \colon X \to Y$，且 $\bar f$ 是
一致连续。\index{一致连续映射的延拓}
\end{theorem}

\begin{proof}
唯一性：两个连续延拓在稠密集 $D$ 上相等，
从而处处相等（集合 $\{g = h\}$ 是闭集：
即闭对角线在 $x\mapsto(g(x), h(x))$ 下的原像）。
存在性：对 $x \in X$，取 $d_n \to x$，$d_n \in D$。
序列 $(f(d_n))$ 是柯西列：给定 $\varepsilon$，由一致
连续性，存在 $\delta$，使 $d(u,v) < \delta \Rightarrow
d(f(u), f(v)) < \varepsilon$，且 $(d_n)$ 为柯西列。令
$\bar f(x) = \lim f(d_n)$；该极限不依赖于
所选序列（交错两个序列即可）。$\bar f$ 延拓 $f$
（取常值序列），并继承连续模：若
若 $d(x, x') < \delta$，用 $D$ 中
距离 $< \frac{\delta - d(x,x')}2$ 的点分别逼近二者，于是 $d(\bar f(x), \bar
f(x')) \leq \varepsilon$（取极限）——$\bar f$ 一致
连续。
\end{proof}

\begin{theorem}[完备化]\label{thm:b3:complete:completion}
每个度量空间 $X$ 都可等距嵌入为
一个完备度量空间 $\hat X$ 的稠密子集，在保持
$X$ 逐点不动的等距意义下唯一：其\emph{完备化}\index{完备化
}。
\end{theorem}

\begin{proof}
\emph{存在性。}设 $\mathcal X$ 为
$X$ 中柯西列的集合，赋以伪距离
\[
D\bigl((x_n), (y_n)\bigr) = \lim_n d(x_n, y_n),
\]
极限存在是因为 $\abs{d(x_n, y_n) - d(x_m, y_m)}
\leq d(x_n, x_m) + d(y_n, y_m)$ 使实数列
为柯西列。令 $\hat X = \mathcal X/{\sim}$，
将 $D$-距离为 $0$ 的序列等同，$D$ 诱导出
一个距离。以常值序列嵌入 $X$，这是等距映射，
其像稠密（柯西列被 $D$-逼近，所用的常值序列
由其自身各项构造：$D\bigl((x_n), (x_k)_{\rm
const}\bigr) = \lim_n d(x_n, x_k) \to 0$ 当 $k \to \infty$ 时，
由柯西性）。$\hat X$ 的完备性：设 $(\xi^k)$ 为
$\hat X$ 中的柯西列；由稠密性取 $x_k \in X$，使 $D(\xi^k, x_k)
\leq 2^{-k}$；于是 $(x_k)$ 是 $X$ 中柯西列（利用
经由 $\xi$ 的三角不等式），确定 $\xi \in \hat
X$ 中的一点，且 $D(\xi^k, \xi) \leq 2^{-k} + D(x_k, \xi) \to 0$（
常值 $x_k$ 到 $(x_j)_j$ 所在类的距离为
$\lim_j d(x_k, x_j)$，当 $k$ 充分大时很小）。

\emph{唯一性}：两个完备化 $\hat X_1, \hat X_2$ 都稠密地
包含 $X$；$X$ 的恒等映射（等距映射）一致连续，
因而可延拓为 $\hat X_1 \to \hat X_2$
（见\cref{thm:b3:complete:extension}），在稠密集上仍为等距映射，
从而处处为等距；反方向同理，于是
两个复合映射固定稠密的 $X$，从而都是恒等映射。
\end{proof}

\begin{theorem}[巴拿赫不动点定理]\label{thm:b3:complete:banach}
设 $X$ 为完备非空空间，$f \colon X \to X$ 为
\emph{压缩映射}：$d(f(x), f(y)) \leq k\,d(x,y)$，其中 $k < 1$。
则 $f$ 有唯一不动点 $x^*$，且每条轨道
都收敛于它，并给出显式收敛速度 $d(x_n, x^*) \leq
\frac{k^n}{1-k}\,d(x_1, x_0)$。\index{巴拿赫不动点定理}
\end{theorem}

\begin{proof}
（本科二年级已证此结论；此处为求自足而重述
其两行论证。）轨道 $x_{n+1} = f(x_n)$ 满足 $d(x_{n+1},
x_n) \leq k^nd(x_1, x_0)$，故为柯西列（等比级数）；
其极限 $x^*$ 是不动点（由 $f$ 的连续性），唯一性因
两个不动点满足 $d \leq k\,d$。收敛速度：对等比
尾项求和。
\end{proof}

\begin{example}[恒等映射的扰动]
\label{ex:b3:complete:perturbident}
设 $g \colon \R^d \to \R^d$ 满足 $k$-利普希茨条件，且 $k < 1$。
则 $\varphi = \mathrm{id} + g$ 为 $\R^d$ 到 $\R^d$ 上的 \emph{同胚}。
单射性，并给出定量的
模估计：
\[
\norm{\varphi(x) - \varphi(y)} \geq \norm{x - y} -
\norm{g(x) - g(y)} \geq (1 - k)\norm{x - y} .
\]
满射性即不动点定理：求解
$\varphi(x) = y$ 即 $x = y - g(x)$，而 $x \mapsto y -
g(x)$ 为 $k$-压缩映射，定义在完备的 $\R^d$ 上 ---
唯一解 $x = \psi(y)$ 对每个 $y$ 均存在。
上式表明逆映射 $\psi$ 满足利普希茨条件，
常数为 $\frac1{1-k}$，因而为同胚，且两个模均有显式
界。这一看似平凡的命题是
反函数定理背后的核心
（见\cref{ch:b3:submanifolds}）：在 $Df$ 可逆的某点附近，
$f$ \emph{恰是}一个可逆线性映射，再加一个
小的利普希茨扰动，本例即可
处理其余部分。此外，它还刻画了数值稳健性：一个系统
若系统受到的扰动小于逆映射的裕度，则仍可解，
且解的移动量至多为扰动量的
$\frac{1}{1-k}$ 倍。
\end{example}

\section{贝尔定理}

\begin{theorem}[贝尔]\label{thm:b3:complete:baire}
在完备度量空间中，稠密开集的可数交
仍稠密。等价地：若 $X = \bigcup_{n}F_n$，其中
每个 $F_n$ 为闭集，则某个 $F_n$ 具有非空
内部。\index{贝尔定理}
\end{theorem}

\begin{proof}
设 $(U_n)$ 为稠密开集，$B_0 = B(x_0, r_0)$ 为任一开
球；求 $\bigcap U_n$ 中位于 $B_0$ 内的一点。归纳构造：
$U_{n}$ 稠密且开，故它与开球 $B_{n-1}$ 相交，
其交为一开集，此开集包含闭球 $\bar B(x_n, r_n)$，
其中 $0 < r_n \leq r_{n-1}/2$，且 $\bar B(x_n, r_n) \subseteq
B_{n-1}\cap U_n$。圆心构成柯西列（$x_m \in
B_n$，当 $m \geq n$ 时，半径 $\to 0$）；极限 $x$ 落在
每个 $\bar B(x_n, r_n)$ 中（由闭性），从而落在每个 $U_n$ 及
$B_0$ 中。关于第二种形式：若每个 $F_n$ 的内部皆空，则
$U_n = X \setminus F_n$ 为开集且稠密，而 $\bigcap U_n$ 中的一点
不属于 $\bigcup F_n = X$，矛盾。
\end{proof}

\begin{remark}\label{rem:b3:complete:meagre}
术语：若集合的闭包
内部为空，则称该集合为\emph{无处稠密}；
若集合是可数个无处稠密集之并，则称其为\emph{贫乏集}\index{贫乏集}（第一纲）。贝尔断言：
\emph{完备度量空间自身不贫乏}，且
贫乏集的补集稠密。“贫乏”是一种与
测度正交的“小”概念（\cref{ch:b3:measure}将
给出满测度的贫乏集），而贝尔型论证可用来
\emph{以丰富性证明存在性}：要
构造一个不具性质 P 的对象，只需证明具性质 P 的对象构成贫乏集。
\end{remark}

\begin{corollary}\label{cor:b3:complete:bairecors}
(a) $\R$ 不可数。(b) $\Q$ 不是可数个
$\R$ 的开子集之交；非空完备度量空间若无孤立点，
则不可数。
\end{corollary}

\begin{proof}
(a) 若 $\R = \bigcup_{x}\{x\}$ 是可数个单点集之并，则某个
单点集将有内部，矛盾。(b) 若 $\Q = \bigcap_n V_n$，其中 $V_n$
为开集（必稠密，因 $\supseteq \Q$），则
$V_n$ 与补集 $\R\setminus\{q\}$（$q \in \Q$）构成一个
可数族稠密开集，其交为空，
与贝尔定理矛盾。若 $X$ 是无孤立点的完备
且可数的空间，则 $X = \bigcup_{x \in X}\{x\}$ 将其表示为
可数个内部为空的闭集之并（因无孤立点）：
再次由贝尔定理得证。
\end{proof}

\begin{theorem}[魏尔斯特拉斯“怪物”的存在性]
\label{thm:b3:complete:nowherediff}
存在 $\intcc01$ 上的连续函数，它在
任何点处都不可导。事实上，$f \in \mathcal C(\intcc01)$ 中
至少在一点具有（有限）导数的全体，在
$(\mathcal C(\intcc01), \norm\cdot_\infty)$ 中为贫乏集。\index{处处
不可导函数}
\end{theorem}

\begin{proof}
对 $n \geq 1$，令
\[
F_n = \Bigl\{f : \exists x \in \intcc01,\
\forall h \neq 0 \text{ 且 } x + h \in \intcc01,\
\abs{f(x+h) - f(x)} \leq n\abs h \Bigr\}.
\]
若 $f$ 在 $x$ 处可导，则 $f \in F_n$ 对某个 $n$ 成立：
商 $\abs{f(x+h)-f(x)}/\abs h$ 在 $\abs h
\leq \delta$ 时有界（由可导性，它趋于
$\abs{f'(x)}$）；商以 $2\norm f_\infty/\delta$ 为界，
对 $\abs h \geq \delta$ 亦然。于是 $\bigcup F_n$ 包含
所有在某点可导的函数，故只需证明
每个 $F_n$ 均为闭集且内部为空。

\emph{闭性}：设 $f_k \to f$ 一致收敛，且 $f_k \in F_n$，
其见证点 $x_k \to x$（由紧性，经子列抽取）。对 $h$
满足 $x + h \in \intcc01$ 者，取 $h_k \to h$，使 $x_k + h_k
\in \intcc01$（例如\ $h_k = h + x - x_k$，截断后）；于是
$\abs{f(x + h) - f(x)} = \lim \abs{f_k(x_k + h_k) - f_k(x_k)}
\leq \lim n\abs{h_k} = n\abs h$，利用一致收敛及
$f$ 在相应各点处的连续性，故 $f \in F_n$。

\emph{内部为空}：给定 $f \in F_n$ 及 $\varepsilon > 0$，
求 $g$，使 $\norm{g - f}_\infty \leq \varepsilon$ 且 $g
\notin F_n$。首先将 $f$ 逼近至误差不超过 $\varepsilon/2$，
用分段仿射函数 $\varphi$（由一致连续性，在
细网格上插值），其斜率以某个 $M$ 为界。再加一个小锯齿：
$g = \varphi + \frac\varepsilon2\,s_N$，其中 $s_N(x)$ 为
周期为 $\frac1N$、振幅为 $1$、斜率为 $\pm 2N$ 的锯齿函数。
对每个 $x$，在其某一侧存在任意小的 $h$，使
锯齿贡献的斜率为 $\pm 2N\cdot\frac\varepsilon2 =
\pm\varepsilon N$，在 $[x, x+h]$ 上，
$g$ 的差商超过 $\varepsilon N - M > n$（当 $N$ 充分大时）。故 $g \notin
F_n$，且在一致距离 $\leq \varepsilon$ 之内（相对 $f$）。
综上，$\bigcup F_n$ 是贫乏集；由贝尔纲定理，其补集
——由处处不可导的函数构成——在
$\mathcal C(\intcc01)$ 中稠密：此类函数\emph{大量
存在}。
\end{proof}

\section{阿尔泽拉–阿斯科利}

全文 $K$ 为\emph{紧}度量空间，且 $\mathcal
C(K) = \mathcal C(K, \R^d)$，范数 $\norm f_\infty = \sup_K
\norm{f(x)}$，是完备空间（连续函数的一致极限
仍是连续的——本科二年级）。

\begin{definition}\label{def:b3:complete:equicontinuous}
称函数族 $\mathcal F \subseteq \mathcal C(K)$ 为
\emph{等度连续}\index{等度连续}，若对每个
$\varepsilon > 0$，存在 $\delta > 0$，使得
\[
d(x, y) < \delta \implies \norm{f(x) - f(y)} < \varepsilon
\quad\text{对\emph{所有} } f \in \mathcal F
\]
（对整个族只取一个 $\delta$——例如任何具有公共利普希茨常数的族，
或具有公共赫尔德模的族），而
\emph{逐点有界}，若 $\sup_{f}\norm{f(x)} < \infty$ 对
每个 $x$ 成立。
\end{definition}

\begin{theorem}[阿尔泽拉–阿斯科利]\label{thm:b3:complete:ascoli}
子集 $\mathcal F \subseteq \mathcal C(K)$ 相对紧
（即其闭包为紧集）当且仅当它等度连续且
逐点有界。特别地，每个等度连续、
逐点有界的\emph{序列}都有
一致收敛的子列。\index{阿尔泽拉–阿斯科利定理}
\end{theorem}

\begin{proof}
（$\Leftarrow$）设 $(f_k)$ 为 $\mathcal F$ 中的序列。
$K$ 为紧度量空间，故是\emph{可分}的：对每个 $n$，有限个
半径为 $\frac1n$ 的球可覆盖 $K$
（见\cref{thm:b3:topology:metriccompact}）；它们的中心构成
可数稠密集 $D = \{x_1, x_2, \dots\}$。由逐点有界性
及波尔查诺–魏尔斯特拉斯定理，依次抽取子列，
使之先在 $x_1$ 处收敛，再在 $x_2$ 处也收敛，依此类推，并取
\emph{对角}子列 $(g_j)$：它在
$D$ 的每一点处都收敛。等度连续性将此升级为一致
柯西性：给定 $\varepsilon$，取定义中的 $\delta$，
以有限个球覆盖 $K$，即 $B(x_i, \delta)$
其中 $x_i \in D$（$i \leq m$），并取 $J$ 充分大，使
$\norm{g_j(x_i) - g_{j'}(x_i)} < \varepsilon$，当 $j, j' \geq
J$，$i \leq m$。对任意 $x \in B(x_i, \delta)$：
\[
\norm{g_j(x) - g_{j'}(x)}
\leq \norm{g_j(x) - g_j(x_i)} + \norm{g_j(x_i) - g_{j'}(x_i)}
+ \norm{g_{j'}(x_i) - g_{j'}(x)} < 3\varepsilon .
\]
于是 $(g_j)$ 为一致柯西列，从而在完备的
$\mathcal C(K)$ 中收敛。于是 $\mathcal F$ 中每个序列都有
收敛子列，所以 $\bar{\mathcal F}$ 是（序列意义下）
紧的，由\cref{thm:b3:topology:metriccompact}可知。

（$\Rightarrow$）若 $\bar{\mathcal F}$ 紧，则逐点
有界性显然（求值映射连续）。对等度连续性，
用有限个球覆盖 $\mathcal F$，
球为 $B(f_i, \varepsilon)$（属于 $\mathcal C(K)$）；每个 $f_i$ 一致连续
（海涅，见\cref{cor:b3:topology:heineborel}），
从而得到公共的 $\delta$，
对 $i \leq m$，任取 $f \in B(f_i, \varepsilon)$ 及满足 $d(x,y)
< \delta$ 时，$\norm{f(x)-f(y)} \leq 2\varepsilon +
\norm{f_i(x)-f_i(y)} < 3\varepsilon$。
\end{proof}

\begin{example}\label{ex:b3:complete:ascoliexamples}
$\mathcal C(\intcc01)$ 中的闭单位球\emph{不是}
紧的（$f_n(x) = x^n$ 没有一致收敛的
子列：其逐点极限不连续），并且
$(x^n)$ 在 $1$ 处不是等度连续的。与此相对，$\{f : \norm
f_\infty \leq 1,\ \operatorname{Lip}(f) \leq 1\}$ 是紧的：
它有界，且为公共利普希茨常数 $1$ 的等度连续族，又是闭的。阿斯科利
解释了\emph{为何}紧性在无穷维空间中失效
（里斯，本科二年级），以及为恢复紧性所需补充的条件：一致
连续模。
\end{example}

\section{斯通–魏尔斯特拉斯}

\begin{lemma}[迪尼]\label{lem:b3:complete:dini}
设 $K$ 为紧集，$(f_n)$ 为一列连续实函数的
\emph{单调}序列，且它\emph{逐点}收敛于一个
连续函数 $f$。则收敛
为一致的。\index{迪尼定理}
\end{lemma}

\begin{proof}
不妨设 $f_n \uparrow f$；令 $g_n = f - f_n \downarrow 0$，
连续。给定 $\varepsilon$，开集 $U_n = \{g_n <
\varepsilon\}$ 递增并覆盖 $K$（由逐点收敛）；
取有限子覆盖：$K = U_{n_0}$，其中 $n_0$ 为某个指标
（递增族），即 $0 \leq g_n < \varepsilon$
处处成立，当 $n \geq n_0$ 时。
\end{proof}

\begin{lemma}\label{lem:b3:complete:sqrt}
存在一列\emph{多项式} $u_n$，使得 $u_n(t) \to
\sqrt t$ 在 $\intcc01$ 上一致收敛。
\end{lemma}

\begin{proof}
令 $u_0 = 0$，$u_{n+1}(t) = u_n(t) + \frac12\bigl(t -
u_n(t)^2\bigr)$，都是多项式。用归纳法证明 $0 \leq u_n(t) \leq
\sqrt t$ 在 $\intcc01$ 上成立：设其对 $n$ 成立，
\[
\sqrt t - u_{n+1}(t) = \bigl(\sqrt t - u_n(t)\bigr)
\Bigl(1 - \tfrac12\bigl(\sqrt t + u_n(t)\bigr)\Bigr) \geq 0,
\]
由于 $\sqrt t + u_n \leq 2$，且 $u_{n+1} \geq u_n \geq 0$。
故 $(u_n(t))$ 单调不减，以 $\sqrt t$ 为上界，因此它
逐点收敛，设极限为 $\ell(t)$，则 $\ell =
\ell + \frac12(t - \ell^2)$，从而 $\ell(t) = \sqrt t$，连续。
由迪尼定理（\cref{lem:b3:complete:dini}）可知此收敛是一致的。
\end{proof}

\begin{theorem}[斯通–魏尔斯特拉斯定理，实值情形]
\label{thm:b3:complete:stoneweierstrass}
设 $K$ 为紧的豪斯多夫空间，$\mathcal A \subseteq
\mathcal C(K, \R)$ 的一个\emph{子代数}（对加法、
乘法和数乘封闭），且\emph{包含常值函数}，
并\emph{分离点}（即对 $x \neq y$，存在 $f \in
\mathcal A$ 使 $f(x) \neq f(y)$）。则 $\mathcal A$ 在
$(\mathcal C(K,\R), \norm\cdot_\infty)$ 中稠密。\index{斯通–魏尔斯特拉斯
定理}
\end{theorem}

\begin{proof}
记 $\bar{\mathcal A}$ 为其闭包，它同样是代数
（在有界集上一致收敛的函数之积仍收敛）。

\emph{第一步：$\bar{\mathcal A}$ 是格}，即关于
$\max$ 与 $\min$ 封闭。由于 $\max(f,g) = \frac{f + g +
\abs{f-g}}2$，$\min$ 同理，只需证 $f \in
\bar{\mathcal A} \Rightarrow \abs f \in \bar{\mathcal A}$：取
$M = \norm f_\infty > 0$，$\abs f = M\sqrt{(f/M)^2}$，
由\cref{lem:b3:complete:sqrt}得多项式列 $u_n$，使
$u_n\bigl((f/M)^2\bigr) \to \abs f/M$ 一致收敛；
而代数中元素的多项式（常数项：常值函数属于代数）
仍属于 $\bar{\mathcal A}$。

\emph{第二步：两点插值。} 对 $x \neq y$ 及 $a,
b \in \R$，存在 $g \in \mathcal A$ 使 $g(x) = a$，$g(y) = b$：
取 $h$ 使其分离 $x, y$，并令 $g = a + (b -
a)\frac{h - h(x)}{h(y) - h(x)}$。

\emph{第三步。} 设 $f \in \mathcal C(K)$，$\varepsilon > 0$。
对每对 $x, y$，取 $g_{x,y} \in \mathcal A$ 使
$g_{x,y}(x) = f(x)$，$g_{x,y}(y) = f(y)$（由第二步；当 $x = y$ 时
取常值函数 $g_{x,x} = f(x)$）。
固定 $x$：对每个 $y$，开集 $V_y = \{g_{x,y} < f +
\varepsilon\}$ 含 $y$。由紧性取出 $y_1, \dots,
y_m$，使 $K = \bigcup V_{y_j}$，且令 $h_x = \min_j g_{x, y_j}
\in \bar{\mathcal A}$（由第一步）满足 $h_x < f +
\varepsilon$ 处处成立，且 $h_x(x) = f(x)$。再令 $x$ 变动：$W_x =
\{h_x > f - \varepsilon\}$ 为开集且含 $x$；取出 $x_1,
\dots, x_l$ 覆盖 $K$，并令 $h = \max_i h_{x_i} \in
\bar{\mathcal A}$ 满足 $f - \varepsilon < h < f +
\varepsilon$，从而 $\norm{f - h}_\infty \leq \varepsilon$。于是
$f \in \bar{\mathcal A}$。
\end{proof}

\begin{corollary}\label{cor:b3:complete:weierstrass}
(a)（\emph{魏尔斯特拉斯}）多项式在 $\mathcal
C(\intcc ab, \R)$ 中稠密；$n$ 元多项式在
$\mathcal C(K, \R)$ 中稠密，其中 $K \subseteq \R^n$ 为紧集。
(b)（\emph{复数情形}）若 $\mathcal A \subseteq \mathcal
C(K, \C)$ 是包含常值函数并分离点的子代数，
且\emph{对共轭封闭}，则它稠密。
(c)（\emph{三角情形}）三角多项式
$\sum_{\abs n \leq N}c_n\eu^{\iu n t}$ 在连续的
$2\pi$ 周期函数空间（赋以
$\norm\cdot_\infty$ 范数）中稠密。\index{魏尔斯特拉斯逼近定理}
\end{corollary}

\begin{proof}
(a) 多项式与常值函数构成代数；坐标
函数分离 $\R^n$ 中的点。
(b) 记 $\frac{f + \bar f}2$，
$\frac{f - \bar f}{2\iu}$ 分别为 $\mathcal A$ 中元素的实部与虚部，它们构成
实代数 $\mathcal A_\R \subseteq \mathcal C(K, \R)$，
含常值函数；它分离点（$f(x) \ne f(y)$ 蕴含
$\operatorname{Re}f$ 或 $\operatorname{Im}f$ 之一能分离这两点）。
对实情形应用定理，再合成。
(c) 将 $2\pi$ 周期函数视为 $\mathcal C(S^1, \C)$
（$S^1 = \R/2\pi\Z$，紧，见
\cref{exo:b3:topology:5}）；由
$\eu^{\iu t}, \eu^{-\iu t}$ 与常值函数生成的代数对共轭封闭，
且分离圆周上的点（$\eu^{\iu t}$
在其上为单射）。应用 (b)。
\end{proof}

\begin{remark}\label{rem:b3:complete:swremark}
三角情形弥补了并极大推广了
本科二年级傅里叶一章留下的缺口：三角多项式在
$(\mathcal C(S^1), \norm\cdot_2)$ 中的稠密性可直接推出，
（在相差归一化常数的意义下 $\norm\cdot_2 \leq \norm\cdot_\infty$），
这将使傅里叶系
在\cref{ch:b3:hilbert}中成为标准正交\emph{基}，
从而最终在完全一般的意义下证明帕塞瓦尔恒等式。
\end{remark}

\section{练习}

\begin{exercise}[$\star$]\label{exo:b3:complete:1}
(a) 证明：$\mathcal C(\intcc01, \R)$ 在
范数 $\norm f_1 = \int_0^1\abs f$ 下\emph{不}完备：
函数 $f_n$ 取值 $0$ 的区域为 $\intcc0{\frac12}$，
取值 $1$ 的区域为 $\intcc{\frac12 + \frac1n}1$，两者之间为仿射函数，
是 $\norm\cdot_1$ 柯西列，却无连续极限。
(b) 证明：若赋范空间中每个绝对收敛的级数都收敛，
则该空间是完备的。\emph{（从
柯西列中抽取子列，使 $\norm{x_{n_{k+1}} -
x_{n_k}} \leq 2^{-k}$。）}
\end{exercise}

\begin{exercise}[$\star$]\label{exo:b3:complete:2}
利用贝尔纲定理：(a) 证明完备赋范空间
不存在可数代数基——由此推出多项式空间
在\emph{任何}范数下都不完备；(b) 证明
若连续函数列 $f_n \colon \R \to \R$ 逐点
收敛于 $f$，则 $f$ 的连续点全体构成的集合稠密。
\emph{（对 (b)：承认或证明 $\Omega_\delta = \{x :
\operatorname{osc}_x f < \delta\}$ 是开集，并证明它
稠密，利用 $F_{N} = \{x: \abs{f_n(x) - f_m(x)} \leq \delta/3\
\forall n,m \geq N\}$，这些闭集覆盖 $\R$；在
任意闭球中进行，以应用贝尔纲定理。）}
\end{exercise}

\begin{exercise}[$\star\star$]\label{exo:b3:complete:3}
(a) 证明 $f(x) = \frac12\bigl(x + \frac ax\bigr)$（$a >
1$）是 $[\sqrt a, +\infty)$ 上的压缩映射，并求出其
不动点——海伦方法。
估计达到 $10^{-12}$ 精度所需的迭代次数，初值取 $x_0 = a$，
当 $a = 2$ 时。
（开普勒方程）设 $0 \leq e < 1$ 及 $m \in \R$，证明
$x = m + e\sin x$ 有唯一解，且该解
关于 $m$ 连续依赖。
\end{exercise}

\begin{exercise}[$\star\star$]\label{exo:b3:complete:4}
(a) 设 $X$ 完备，$f \colon X \to X$ 使得某个
迭代 $f^p$ 是压缩映射。证明 $f$ 有唯一
不动点。应用：积分算子 $T$ 作用于
$\mathcal C(\intcc0a)$ 上，$Tf(x) = \int_0^x f$，满足
$\norm{T^p} \leq a^p/p!$——由此推出
$u = g + \lambda Tu$ 对每个 $\lambda$ 都可解。
(b)（埃德尔斯坦）设 $K$ 是\emph{紧}的，$f\colon K \to K$
且 $d(f(x), f(y)) < d(x, y)$，对 $x \neq y$ 成立。证明 $f$
有唯一不动点，但压缩系数可能
丢失：在 $X = [1, +\infty)$（完备但非紧）上，$f(x) = x
+ \frac1x$ 没有不动点，尽管距离严格
递减。
\end{exercise}

\begin{exercise}[$\star\star$]\label{exo:b3:complete:5}
(a) 两个取值于豪斯多夫空间的连续映射，若二者在一个
稠密子集上一致，则二者处处相等；本章何处用到了
这一结论？
(b) 设 $D \subseteq X$ 为稠密子集，$f \colon D \to Y$ 是
到完备空间 $Y$ 的稠密子集上的等距双射，且
$X$ 完备。证明 $f$ 可延拓为等距双射
$X \to Y$。由此再次推出完备化的唯一性。
\end{exercise}

\begin{exercise}[$\star\star$]\label{exo:b3:complete:6}
下列哪些函数族是等度连续的、逐点
有界的、在 $\mathcal C(\intcc01)$ 中相对紧的？
\[
\{x \mapsto \sin(nx)\}_{n},\qquad
\{x \mapsto x^n\}_n,
\]
\[
\Bigl\{f \in \mathcal C^1 : \norm f_\infty \leq 1,\
\norm{f'}_\infty \leq 1\Bigr\},\qquad
\{f : \operatorname{Lip}(f)\leq 1,\ f(0) = 0\}.
\]
每个答案都须借助阿尔泽拉–阿斯科利定理或一个反例函数列加以论证。
\end{exercise}

\begin{exercise}[$\star\star\star$]\label{exo:b3:complete:7}
（紧积分算子）设 $k \in \mathcal C(\intcc01^2)$
且对 $f \in \mathcal C(\intcc01)$，记 $Tf(x) =
\int_0^1k(x,y)f(y)\,\dd y$.
(a) 证明 $T$ 把 $\mathcal C(\intcc01)$ 的单位球
映为等度连续、一致有界的集合；由此断定 $T$
是\emph{紧算子}：有界集的像
均为相对紧集。
(b) 由此推出：若 $(f_n)$ 有界，则 $(Tf_n)$ 有一致收敛
的子列，且 $T$ 不可能是具有
连续逆的双射。\emph{（单位球的像将是
$0$ 在 $\mathcal C(\intcc01)$ 中的一个紧邻域；
这被本科二年级的里斯定理排除。）}
\end{exercise}

\begin{exercise}[$\star\star$]\label{exo:b3:complete:8}
对连续函数 $f_n \colon \intcc01 \to \R$，证明或否定：
(a) $f_n \downarrow 0$ 逐点收敛 $\Rightarrow$ 一致收敛（狄尼
——请重新证明）；(b) 去掉单调性；(c) 保留单调性，
但 $f$ 不连续；(d) 保留单调性及连续极限，
但定义在 $\intoo01$ 上。
\end{exercise}

\begin{exercise}[$\star\star$]\label{exo:b3:complete:9}
(a)（矩唯一决定）设 $f \in \mathcal C(\intcc01)$，满足
$\int_0^1 f(x)\,x^n\,\dd x = 0$ 对所有 $n \in \N$ 成立，证明 $f =
0$。\emph{（用多项式一致逼近 $f$，并计算
$\int f^2$。）}
(b) 证明偶多项式在 $\mathcal
C(\intcc01)$ 中稠密，但在 $\mathcal C(\intcc{-1}1)$ 中\emph{不}稠密；
斯通–魏尔斯特拉斯定理的哪一条假设不成立？
(c) 仅由 $x \mapsto \eu^{\iu x}$ 生成的代数
（不含 $\eu^{-\iu x}$）是否在 $\mathcal C(S^1, \C)$ 中稠密？
\emph{（考虑 $\int_0^{2\pi}f(t)\,\eu^{\iu t}\dd t$。）}
\end{exercise}

\begin{exercise}[$\star\star\star$]\label{exo:b3:complete:10}
（一致有界性，度量版本）设 $X$ 为完备
度量空间，$\mathcal F \subseteq \mathcal C(X, \R)$ 是一个
\emph{逐点}有界的函数族：$\sup_{f \in \mathcal
F}\abs{f(x)} < \infty$（对每个 $x$）。证明存在一个
非空开集 $U \subseteq X$，使 $\mathcal F$ 在其上
\emph{一致}有界：$\sup_{f}\sup_U \abs f < \infty$。
\emph{（考虑 $F_n = \{x : \abs{f(x)} \leq n\ \forall f\}$。）}
这正是巴拿赫–斯坦豪斯定理背后的动力，
该定理在\cref{ch:b3:banach}中给出。
\end{exercise}

\begin{exercise}[$\star\star$]\label{exo:b3:complete:11}
（$\mathcal C^1$ 需要其自身的范数）在 $E = \mathcal
C^1(\intcc01, \R)$ 上考虑 $\norm f_{\mathcal C^1} = \norm
f_\infty + \norm{f'}_\infty$。
(a) 证明 $(E, \norm\cdot_{\mathcal C^1})$ 是完备的，
\emph{（$\mathcal C^1$-柯西列满足 $f_n \to f$ 且
$f_n' \to g$ 一致收敛；通过取极限确定 $g = f'$，
即在 $f_n(x) = f_n(0) + \int_0^xf_n'$ 中取极限）}。
(b) 证明 $(E, \norm\cdot_\infty)$ \emph{不}
是完备的：给出一列 $\mathcal C^1$ 函数的一致极限，
该极限不可导（例如
$\abs{x - \tfrac12}$ 的光滑逼近）。
(c) 借助 (a)、(b) 以及开映射定理一类的思想
--- 或直接 --- 推出不存在常数 $C$ 满足
$\norm{f'}_\infty \leq C\,\norm f_\infty$ 在 $E$ 上：给出
一个见证该结论的序列。求导算子无界，这正是
\cref{thm:b3:complete:nowherediff}背后的断崖。
\end{exercise}

\begin{exercise}[$\star\star\star$]\label{exo:b3:complete:12}
（克罗夫特引理）设 $f\colon\intoo0\infty\to\R$ 连续，
并假设对每个 $x > 0$，$f(nx) \to 0$，
其中整数 $n \to \infty$。证明 $f(t) \to 0$，当 $t \to
+\infty$ 时。\emph{（固定 $\varepsilon > 0$；集合
\[
F_N = \{x > 0 : \abs{f(nx)} \leq \varepsilon\ \ \forall n
\geq N\}
\]
这些集合是闭集，且覆盖 $\intoo0\infty$；由贝尔纲定理，在某区间
$\intcc ab$ 中可得 $N$ 及子区间 $\intcc{a'}{b'}
\subseteq F_N$；于是伸缩后的区间 $\bigl[na', nb'\bigr]$，$n
\geq N$ 的这些区间覆盖 $+\infty$ 的一个完整邻域，只要
$n(b' - a') \geq a'$。）}假设“对每个
$x$”（而非仅有理数 $x$）这一条件在何处被用到？
\end{exercise}

\section{问题：佩亚诺存在性定理}

\begin{problem}[{周末问题 --- 无利普希茨条件时
$x' = f(t, x)$ 解的存在性}]\label{pb:b3:complete:1}
柯西–利普希茨定理（本科二年级；在\cref{ch:b3:ode}中重新证明）
要求 $f$ 关于 $x$ 满足利普希茨条件。佩亚诺（1890）：\emph{$f$ 的连续性
本身已足以保证解的存在性} --- 但不保证唯一性。
下面借助欧拉折线与阿斯科利定理证明之。设 $f \colon R
\to \R^d$ 在矩形 $R = \intcc{t_0 - a}{t_0
+ a}\times \bar B(x_0, b) \subseteq \R\times\R^d$ 上连续，$M =
\sup_R\norm f$，且
\[
T = \min\bigl(a,\ b/M\bigr)
\qquad (\text{若 } M = 0 \text{ 则问题平凡}).
\]

\medskip
\noindent\textbf{第一部分 --- 欧拉折线。}对 $n \geq 1$，
将 $[t_0, t_0 + T]$ 按 $t_k = t_0 + kT/n$（$0 \leq k
\leq n$），并分段仿射地定义 $\varphi_n$：
$\varphi_n(t_0) = x_0$，且在 $[t_k, t_{k+1}]$ 上，
\[
\varphi_n(t) = \varphi_n(t_k) + (t -
t_k)\,f\bigl(t_k, \varphi_n(t_k)\bigr).
\]
\begin{enumerate}
  \item 用归纳法证明 $\varphi_n$ 是良定义的，
        且 $\norm{\varphi_n(t) - x_0} \leq M(t - t_0) \leq b$
        在 $[t_0, t_0 + T]$ 上 --- 因此取值点仍留在
        $R$ 内。\emph{（此处用到了 $T \leq b/M$。）}
  \item 证明每个 $\varphi_n$ 都是 $M$ 利普希茨的。
  \item 由阿尔泽拉–阿斯科利
        （\cref{thm:b3:complete:ascoli}）可知存在子列
        $(\varphi_{n_j})$ 在 $[t_0, t_0 +
        T]$ 上一致收敛于某个 $\varphi$，它亦为 $M$ 利普希茨的，
        且 $\varphi(t_0) = x_0$。
\end{enumerate}

\medskip
\noindent\textbf{第二部分 --- 极限函数是方程的解。}
定义\emph{偏差} $\Delta_n(t) = \varphi_n'(t) -
f\bigl(t, \varphi_n(t)\bigr)$ 在非分点处。
\begin{enumerate}[resume]
  \item 证明 $f$ 在 $R$ 上一致连续，并由此推出：
        对每个 $\varepsilon > 0$，存在 $n_0$，使得
        当 $n \geq n_0$ 且每个非分点 $t$ 时，
        $\norm{\Delta_n(t)} \leq \varepsilon$。\emph{（在 $(t_k,
        t_{k+1})$ 上 $\varphi_n'(t) = f(t_k, \varphi_n(t_k))$，
        且 $(t, \varphi_n(t))$ 与参照点的距离至多为 $(1 + M)\,T/n$，
        参照点为 $(t_k, \varphi_n(t_k))$。）}
  \item 证明积分形式：对所有 $t$，
        \[
        \varphi_n(t) = x_0 + \int_{t_0}^{t}
        f\bigl(s, \varphi_n(s)\bigr)\dd s +
        \int_{t_0}^t \Delta_n(s)\,\dd s,
        \]
        其中中间的被积函数分段连续。
  \item 沿 $(n_j)$ 取极限，证明
        $f(s, \varphi_{n_j}(s)) \to f(s, \varphi(s))$
        \emph{一致地}（再次利用 $f$ 的一致连续性），并
        得出
        \[
        \varphi(t) = x_0 + \int_{t_0}^t
        f\bigl(s, \varphi(s)\bigr)\dd s .
        \]
  \item 由此推出 $\varphi$ 为 $\mathcal C^1$ 函数，在 $[t_0, t_0 +
        T]$ 上，且满足 $x' = f(t, x)$，$x(t_0) = x_0$；
        将构造推广到 $[t_0 - T, t_0]$（时间反演）。
        \emph{此即佩亚诺定理。}
\end{enumerate}

\medskip
\noindent\textbf{第三部分 --- 唯一性确实不成立。}
考虑 $x' = 2\sqrt{\abs x}$，$x(0) = 0$，定义在 $\R$ 上。
\begin{enumerate}[resume]
  \item 验证 $f(x) = 2\sqrt{\abs x}$ 连续，
        但在 $0$ 的任何邻域上都不满足利普希茨条件。
  \item 验证 $x \equiv 0$ 以及对每个 $c \geq 0$，
        \[
        x_c(t) = \begin{cases} 0 & t \leq c,\\ (t - c)^2 & t
        \geq c,\end{cases}
        \]
        都是过 $(0,0)$ 的解：共得一个连续统的
        互异解。
  \item 对于这个 $f$，皮卡迭代论证（巴拿赫不动点）
        在何处失效？
\end{enumerate}

\medskip
\noindent\textbf{第四部分 --- 方法的极限。}
\begin{enumerate}[resume]
  \item 证明皮亚诺定理在无穷维中不成立：
        这里可承认（或直接接受）经典的
        迪厄多内在趋于零的数列空间 $c_0$ 中的例子；
        转而证明在那里失效的有限维要素：
        即 $c_0$ 的闭单位球（上确界范数）
        不是紧的 --- 试给出一个有界数列，
        它没有收敛子列，并说明
        第一部分中哪一步在此失效。
  \item 总结：哪些假设给出存在性？哪些给出
        存在性与唯一性？请精确叙述下面两个定理
        （皮亚诺；柯西--利普希茨），并列对照。
\end{enumerate}

\medskip
\noindent\textbf{第五部分 --- 奥斯古德：超出
利普希茨条件的唯一性。} 设 $\omega\colon\intoo0\infty\to\intoo0\infty$ 为
连续、非减，且满足
\[
\int_0^1\frac{\dd r}{\omega(r)} = +\infty
\qquad(\text{发散于 } 0),
\]
并设 $f$ 满足 $\norm{f(t, x) - f(t, y)} \leq
\omega\bigl(\norm{x - y}\bigr)$ 在 $R$ 上。
\begin{enumerate}[resume]
  \item 验证 $\omega(r) = Lr$ 满足条件（即利普希茨条件），
        而 $\omega(r) = r\log\frac1r$（由连续性延拓，对
        小的 $r$）也满足条件，尽管它 \emph{不是} $O(r)$，
        以及 $\omega(r) = 2\sqrt r$ 不满足条件。计算
        各情形下的积分。
  \item 设 $x_1, x_2$ 为方程在 $[t_0, t_0 + T]$ 上的解，
        满足 $x_1(t_0) = x_2(t_0)$，并记 $\delta(t) =
        \norm{x_1(t) - x_2(t)}$。仅由积分形式
        本身证明：当 $t_0 \leq s \leq t$ 时，
        \[
        \delta(t) \leq \delta(s) +
        \int_s^t\omega\bigl(\delta(v)\bigr)\dd v .
        \]
  \item （奥斯古德定理）设 $\delta(t_1) > 0$，其中
        $t_1$ 为某个时刻，并令 $\tau = \sup\{t \leq t_1 : \delta(t) =
        0\}$。对 $s \in \intoo\tau{t_1}$，令 $u(t) = \delta(s)
        + \int_s^t\omega(\delta(v))\,\dd v$。证明 $\delta \leq
        u$，$u' = \omega(\delta) \leq \omega(u)$，并推出
        \[
        \int_{\delta(s)}^{u(t_1)}\frac{\dd r}{\omega(r)}
        \leq t_1 - s \leq t_1 - \tau .
        \]
        令 $s \downarrow \tau$，并由
        积分的发散性导出矛盾。于是得出结论：
        \emph{经过同一初始条件的解
        必然重合} --- 奥斯古德条件下的唯一性。
  \item 推出以下结论：柯西--利普希茨唯一性是
        $\omega(r) = Lr$ 的情形；方程 $x' =
        x\log\frac1{\abs x}$（在 $0$ 处延拓为 $0$）有唯一
        解，尽管其右端在
        $0$ 处不满足利普希茨条件；而对 $x' = 2\sqrt{\abs x}$，
        $\int_0\frac{\dd r}{2\sqrt r}$ 的收敛性恰好使
        解能在有限时间内离开 $0$ --- 请将
        积分值 $\int_0^{h^2}\frac{\dd r}{2\sqrt r} = h$
        与 $x_c$ 的离开行为对照。
\end{enumerate}

\medskip
\noindent\textbf{第六部分 --- 速率、格式与漏斗。}
\begin{enumerate}[resume]
  \item （积分形式的格朗沃尔引理）设 $e, \eta \geq 0$ 在
        $[t_0, t_0+T]$ 上连续，且 $e(t) \leq \int_{t_0}
        ^t\bigl(L\,e(s) + \eta(s)\bigr)\dd s$ 对所有 $t$ 成立。证明
        \[
        e(t) \leq \frac{\sup\eta}{L}\bigl(\eu^{L(t - t_0)} -
        1\bigr)
        \]
        \emph{（令 $G(t) = \int_{t_0}^t(Le + \eta)$，注意 $G'
        \leq LG + \sup\eta$，并对
        $\eu^{-Lt}G(t)$ 求导）}。
  \item （欧拉方法具有收敛速率）现设 $f$ 满足：
        $L$-利普希茨（关于 $x$），$L'$-利普希茨（关于 $t$），在 $R$ 上。
        结合第 4 题中的误差界（精确化为
        $\norm{\Delta_n} \leq (L' + LM)T/n$），并与
        第 17 题相结合，证明
        \[
        \sup_{[t_0, t_0+T]}\norm{\varphi_n - \varphi}
        \;\leq\; \frac{(L' + LM)\,T}{n}\cdot
        \frac{\eu^{LT} - 1}{L} = O\Bigl(\frac1n\Bigr),
        \]
        其中 $\varphi$ 即\emph{唯一}解：在利普希茨数据下，
        整个数列收敛，并给出显式
        速率 --- 无需子列。为何唯一性
        能将子列收敛升级为全序列收敛，
        即使不做此计算？
  \item （格式决定结果）对于 $x' = 2\sqrt{\abs x}$，$x(0) =
        0$：证明每条欧拉折线恒为零，
        因此格式收敛到解 $x \equiv 0$；
        但若从 $x(0) = \varepsilon > 0$ 出发，则
        （先令 $n \to \infty$，再令 $\varepsilon \to 0$）收敛到 $t
        \mapsto t^2$，是经过原点的\emph{另一个}解。
        非唯一性以数值格式对扰动的敏感性的形式重新出现。
        numerical scheme to perturbations.
  \item 证明：记 $\mathcal S$ 为 \emph{全部}
        $x' = f(t,x)$、$x(t_0) = x_0$ 在 $[t_0, t_0+T]$ 上的解所成的集合，
        （取值于 $\bar B(x_0, b)$）非空（第二部分），
        一致 $M$-利普希茨，且在
        $\bigl(\mathcal C([t_0, t_0+T], \R^d),
        \norm\cdot_\infty\bigr)$ 中闭；由阿斯科利定理得
        $\mathcal S$ 是\emph{紧}的。（克内泽尔定理还断言
        $\mathcal S$ 是连通的，
        但我们不予证明。）
  \item 在该例上验证克内泽尔现象：对 $x' =
        2\sqrt{\abs x}$，$x(0) = 0$，在 $[0, 1]$ 上证明
        $\mathcal S = \{x_c : c \in \intcc0\infty\}$，其中
        $x_\infty \equiv 0$ \emph{（对任一解，令 $c =
        \sup\{t : x(t) = 0\}$，并对 $(\sqrt x)' = 1$ 在
        $\{x > 0\}$ 上积分）}，证明 $c \mapsto x_c$ 自
        $\intcc0\infty$（单点紧化，即
        令 $c = \infty$ 作为极限补上）到 $\mathcal C([0,1])$ 连续，
        并得出 $\mathcal S$ 确实是紧的且
        连通的 --- 一个线段形的漏斗。
  \item （可达集）由第 20 题推出：对每个
        固定的 $t$，\emph{可达集} $\mathcal S(t) =
        \{x(t) : x \in \mathcal S\}$ 是紧的；对
        第 21 题的例子计算它，并验证它也是
        连通的：$\mathcal S(t) = \intcc0{t^2}$ --- 每个
        中间状态都被某个解取到。
\end{enumerate}

\medskip
\noindent\textbf{第七部分 --- 补充：依赖性、
最优性与一种手算格式。}
\begin{enumerate}[resume]
  \item （连续依赖性）设 $f$ 满足 $L$-利普希茨条件，
        即关于 $x$ 在 $R$ 上成立；又设 $x, y$ 为两个解，
        初值分别为 $x_0, y_0$（在 $t_0$ 处给定）。将
        第 17 题的证明改写，使其适用于不等式 $e(t) \leq \norm{x_0
        - y_0} + \int_{t_0}^tL\,e(s)\dd s$，证明
        \[
        \norm{x(t) - y(t)} \leq \norm{x_0 - y_0}\,
        \eu^{L(t - t_0)},
        \]
        并在 $x' = Lx$ 上验证该界可以取到：
        格朗沃尔不等式是最优的。再由此推出唯一性（$x_0 =
        y_0$），以及流映射 $x_0 \mapsto x(t; x_0)$
        是利普希茨映射，常数为 $\eu^{LT}$，在其
        有定义之处均成立。
  \item （奥斯古德条件是最优的）反之，设 $\omega$ 为
        连续、非减、在 $\intoo0\infty$ 上为正的函数，
        且 $\omega(r) \to 0$，当 $r \to 0^+$ 时，
        但
        \[
        \int_0^1\frac{\dd r}{\omega(r)} < +\infty,
        \]
        并延拓 $f(x) = \omega(\abs x)$，取 $f(0) = 0$。
        证明 $\Omega(x) = \int_0^x\frac{\dd r}{\omega(r)}$
        是从 $\intoc0{1}$ 到
        $\intoc0{\Omega(1)}$ 上的递增双射，其逆 $g$ 满足 $g'
        = \omega(g)$，且 $g(0^+) = 0$，并证明 $g$，
        以 $0$ 延拓至 $t \leq 0$ 时，是
        $\mathcal C^1$ 类的 $x' = f(x)$ 的解，过 $(0, 0)$，且异于 $x \equiv
        0$。\emph{（对 $g'(0) = 0$ 的情形，将 $\frac{g(t)}t$ 以
        $\omega(g(t))$）}。由此可知，
        第 15 题中的发散条件并非仅是便利的假设，而是唯一性的
        精确边界；由
        $\omega(r) = 2\sqrt r$，$\Omega(x) = \sqrt x$ 即可重新得到第三部分的结果。
  \item （欧拉手算）对 $x' = x$，$x(0) = 1$，在
        $[0, 1]$ 上：证明欧拉折线满足
        $\varphi_n(1) = \bigl(1 + \frac1n\bigr)^n$。证明下列
        展开式
        \[
        \Bigl(1 + \frac1n\Bigr)^{n}
        = \eu\Bigl(1 - \frac1{2n} + O\Bigl(\frac1{n^2}\Bigr)
        \Bigr),
        \]
        于是 $t = 1$ 处的误差为 $\frac{\eu}{2n} +
        O(n^{-2})$：第 18 题中的 $O(\frac1n)$ 阶，并给出
        精确常数。对 $n = 10$ 做数值检验：
        $1.1^{10} = 2.59374$，对照 $\eu \approx 2.71828$，
        误差为 $0.12454$，可与 $\frac{\eu}{20}
        \approx 0.13591$。
\end{enumerate}
\end{problem}
